\subsection{An error bound for the scalar limit}

Suppose that $f^{1/n}$ belongs to the space $\mathcal D_h^0$ of \cref{subsec:bounded-root-approximation}. The uniform tail condition of \cref{prop:root-uniform-tails} then gives a rate for the convergence $c_\lambda\to b_h(f)$ in \cref{cor:individual-positive-threshold-exponential}. Split the root $p=f^{1/n}$ at a height $K$. On $\{p\leq K\}$, the error $p(1-\mathrm{e}^{\lambda v/n})$ of the damped root has norm $O(K\lambda)$ in $\mathcal D_h$, uniformly over normalized potentials $v$, while on $\{p>K\}$ it is bounded by the tail of $p$, which is small in that norm. Optimizing over $K$ gives a modulus valid for every normalized potential.

\begin{proposition}\label{prop:exponential-root-tail-modulus}
Let $0\leq p\in\mathcal D_h^0$ be nonzero, and put $f=p^n$, $R=\rho_h(p)$ and $b=R^{-n}=b_h(f)$. There is a constant $C_h$, depending only on $h$ and $\mathrm{d}V$, such that the function
\[
 \eta_p(\lambda)=\inf_{K>0}
 \left(\|p\mathbf1_{\{p>K\}}\|_{\mathcal D_h}+C_hK\lambda\right)
\]
tends to zero as $\lambda\downarrow0$ and satisfies
\begin{equation}\label{eq:uniform-root-damping}
 \big\|p(1-\mathrm{e}^{\lambda v/n})\big\|_{\mathcal D_h}
 \leq\eta_p(\lambda)
 \qquad\bigl(v\in\PSH(X,h),\ \sup_Xv=0\bigr).
\end{equation}
Let $u_\lambda$, $v_\lambda$ and $c_\lambda$ be the greatest exponential solutions and their normalizations from \cref{cor:individual-positive-threshold-exponential}. If $\eta_p(\lambda)<R$, then
\[
 b\leq c_\lambda\leq\bigl(R-\eta_p(\lambda)\bigr)^{-n}
\]
and
\begin{align*}
 \big\|\mathrm{MA}^{\mathrm{ext}}_h(v_\lambda)
              -bf\,\mathrm{d}V\big\|_{\mathrm{TV}}
 &\leq(c_\lambda-b)\|p\|_{L^n}^{n}
       +nb\|p\|_{L^n}^{n-1}\eta_p(\lambda).
\end{align*}
If $p\in L^r$ for some finite $r>n$, the same bounds hold with $\eta_p(\lambda)$ replaced by $C_{h,r}\|p\|_{L^r}\lambda$, where $C_{h,r}$ depends only on $h$, $\mathrm{d}V$ and $r$. In particular, for each such $p$ the scalar and total-variation errors are $O(\lambda)$.
\end{proposition}
\Needspace{10\baselineskip}
\begin{proof}
By \cref{lem:uniform-exponential}, normalized $h$-plurisubharmonic functions are uniformly bounded in every $L^s$ with $s<+\infty$. By \eqref{eq:root-subcritical-bound} with $r=2n$, their norms $\|v\|_{\mathcal D_h}$ are therefore uniformly bounded; let $C_h$ be $1/n$ times this bound. Let $v$ be normalized and $K>0$. On $\{p\leq K\}$ we have
\[
 0\leq p(1-\mathrm{e}^{\lambda v/n})
     \leq K\lambda(-v)/n,
\]
since $1-\mathrm{e}^{-x}\leq x$ for $x\geq0$, while on $\{p>K\}$ we use the bound $p(1-\mathrm{e}^{\lambda v/n})\leq p$. The monotonicity of $\|\cdot\|_{\mathcal D_h}$ and the triangle inequality give
\[
 \big\|p(1-\mathrm{e}^{\lambda v/n})\big\|_{\mathcal D_h}
 \leq\|p\mathbf1_{\{p>K\}}\|_{\mathcal D_h}+C_hK\lambda,
\]
and taking the infimum over $K$ proves \eqref{eq:uniform-root-damping}. By \cref{prop:root-uniform-tails}, and since the $L^n$ tails of $p$ tend to zero, $\|p\mathbf1_{\{p>K\}}\|_{\mathcal D_h}\to0$ as $K\to\infty$. Choosing $K$ first and then letting $\lambda\downarrow0$ shows that $\eta_p(\lambda)\to0$.

Now set $q_\lambda=p\,\mathrm{e}^{\lambda v_\lambda/n}$, so that $q_\lambda^n=g_\lambda=\mathrm{e}^{\lambda v_\lambda}f$, with $g_\lambda$ as in the proof of \cref{cor:individual-positive-threshold-exponential}. Since $0\leq q_\lambda\leq p$ and $\rho_h(p-q_\lambda)\leq\eta_p(\lambda)$ by \eqref{eq:uniform-root-damping}, the triangle inequality and monotonicity give
\[
 R-\eta_p(\lambda)\leq\rho_h(q_\lambda)\leq R.
\]
As in the proof of \cref{cor:individual-positive-threshold-exponential}, $b\leq c_\lambda\leq b_h(g_\lambda)=\rho_h(q_\lambda)^{-n}$, which yields the scalar bound. For the measures, the factorization of $p^n-q_\lambda^n$ gives $0\leq p^n-q_\lambda^n\leq np^{n-1}(p-q_\lambda)$, and H\"older's inequality implies
\[
 \|f-g_\lambda\|_{L^1}
 \leq n\|p\|_{L^n}^{n-1}\|p-q_\lambda\|_{L^n}
 \leq n\|p\|_{L^n}^{n-1}\eta_p(\lambda).
\]
Since $\mathrm{MA}^{\mathrm{ext}}_h(v_\lambda)=c_\lambda g_\lambda\,\mathrm{d}V$, the identity $c_\lambda g_\lambda-bf=(c_\lambda-b)g_\lambda-b(f-g_\lambda)$, together with $\|g_\lambda\|_{L^1}\leq\|f\|_{L^1}=\|p\|_{L^n}^n$, proves the total-variation bound.

Finally, suppose $p\in L^r$ with $n<r<+\infty$, and choose $n<s<r$. By \eqref{eq:root-subcritical-bound} with exponent $s$, and since $\|\cdot\|_{L^n}\leq\|\cdot\|_{L^s}$, we have $\|\cdot\|_{\mathcal D_h}\leq(1+C_s)\|\cdot\|_{L^s}$. H\"older's inequality and \cref{lem:uniform-exponential} then give, for every normalized $v$,
\[
 \|p(-v)\|_{\mathcal D_h}
 \leq(1+C_s)\|p(-v)\|_{L^s}
 \leq(1+C_s)\|p\|_{L^r}\|v\|_{L^{rs/(r-s)}}
 \leq nC_{h,r}\|p\|_{L^r}.
\]
Since $0\leq1-\mathrm{e}^{\lambda v/n}\leq\lambda(-v)/n$, the left side of \eqref{eq:uniform-root-damping} is therefore at most $C_{h,r}\|p\|_{L^r}\lambda$, and the scalar and total-variation estimates can be repeated with this bound in place of $\eta_p(\lambda)$. For $\lambda$ so small that $C_{h,r}\|p\|_{L^r}\lambda\leq R/2$, the function $t\mapsto(R-t)^{-n}$ has bounded derivative on $[0,R/2]$, so the scalar error $c_\lambda-b$, and hence the total-variation error, is $O(\lambda)$.
\end{proof}

For a general $p\in\mathcal D_h^0$ the modulus $\eta_p(\lambda)$ tends to zero, but it need not be $O(\lambda)$. The rescaled scalar error $\lambda^{-1}\log(c_\lambda/b)$ is the loss in the lower bound $v_\lambda\geq v-\lambda^{-1}\log(c_\lambda/b)$ for normalized threshold subsolutions $v$, and the scalar bound of the proposition gives
\[
 0\leq\lambda^{-1}\log(c_\lambda/b)
 \leq-\frac n\lambda\log\bigl(1-\eta_p(\lambda)/R\bigr),
\]
which stays bounded as $\lambda\downarrow0$ when $\eta_p(\lambda)=O(\lambda)$, in particular for $p\in L^r$ with $r>n$; root tails alone do not give such a bound. Indeed, on the products of \cref{sec:zero-lower-mass-product-equations} there are densities whose roots lie in $\mathcal D_h^0$ and whose rescaled errors $\lambda^{-1}\log(c_\lambda/b_h(f))$ tend to $+\infty$ (\cref{sec:nonlp-selection}); for their roots $\eta_p(\lambda)/\lambda$ is unbounded. Even when a uniform lower barrier for $v_\lambda$ is available, passing to the limit in capacity requires the singularity-type or lower-mass hypotheses of a continuity theorem. The proposition estimates only scalars and measures; when the lower mass vanishes, it does not give an exact equation at the threshold.
